\section{Sharp Accuracy Bounds from Simultaneous Payoff Contrasts}\label{app:r18closure}
Fix a finite nonempty set $I$ of computed procedures, all evaluated on the
same declared economic population. For each supplied pair $(a,b)$, let
$[L_{ab},U_{ab}]$ be a finite interval for $V_a-V_b$. All intervals are
assumed to hold on one event $\mathcal E$ with probability at least
$1-\alpha$. No independence across intervals is required. Define
\[
 \mathcal U=\{v\in\R^I:L_{ab}\leq v_a-v_b\leq U_{ab}
                    \text{ for every supplied pair}\}.
\]
The graph has an edge $b\to a$ of length $U_{ab}$ and an edge $a\to b$
of length $-L_{ab}$ for every supplied pair. It is connected in both
directions in the application, because every candidate has a comparison
with the reference. The convention for the empty path is length zero.

\begin{proposition}[Exact catalogue closure]\label{prop:r18closure}
Suppose the graph is strongly connected and contains no negative cycle.
Let $D_{ij}$ be its shortest-path lengths. Then $\mathcal U$ is nonempty,
$D_{ii}=0$, and
\begin{align}
 \sup_{v\in\mathcal U}(v_j-v_i)&=D_{ij},\label{eq:r18sharpcontrast}\\
 \inf_{v\in\mathcal U}\{\max_j v_j-v_i\}
     &=\max_j(-D_{ji}),\label{eq:r18sharplower}\\
 \sup_{v\in\mathcal U}\{\max_j v_j-v_i\}
     &=\max_j D_{ij}.\label{eq:r18sharpupper}
\end{align}
All three extrema are attained. Normalizing the reference value to zero
does not change them. On $\mathcal E$, these bounds hold simultaneously
for every candidate and every pair, including pairs not directly compared
in the original report.
\end{proposition}
\begin{proof}
Every edge $a\to b$ of length $w$ imposes $v_b-v_a\leq w$.
Summing along a path and then minimizing yields
$v_j-v_i\leq D_{ij}$ for every $v\in\mathcal U$.
There are finite shortest paths because the graph is finite, strongly
connected, and has no negative cycle. They satisfy the triangle inequality
$D_{ib}\leq D_{ia}+D_{ab}$ and $D_{ii}=0$.
For fixed $i$, the vector $v_k^+=D_{ik}$ satisfies each original edge
constraint: if $a\to b$ has length $w$, then
$D_{ib}\leq D_{ia}+w$. It therefore belongs to $\mathcal U$ and
attains $v_j^+-v_i^+=D_{ij}$ simultaneously for all $j$. This proves
\eqref{eq:r18sharpcontrast} and \eqref{eq:r18sharpupper}.

For the lower bound, reversing the pairwise inequality gives
$v_j-v_i\geq-D_{ji}$. The vector $v_k^-=-D_{ki}$ is also feasible,
since the same edge implies $D_{ai}\leq w+D_{bi}$ and hence
$v_b^--v_a^-=D_{ai}-D_{bi}\leq w$. It simultaneously attains
$v_j^--v_i^-=-D_{ji}$, proving \eqref{eq:r18sharplower}.
Both extrema include $j=i$, so both regret bounds are nonnegative.
Subtracting the reference coordinate leaves every difference unchanged.
Finally, on $\mathcal E$ the true payoff vector belongs to $\mathcal U$;
all deterministic conclusions therefore hold on that same event.
\end{proof}

Sharpness here is relative to the interval-defined uncertainty set. It does
not assert that every compatible vector can be generated by the underlying
capital model, or that the original statistical intervals are optimal.
A negative cycle makes $\mathcal U$ empty. The implementation stops with
an inconsistency diagnostic; it does not delete an unfavorable interval,
truncate a negative cycle, or report an accuracy certificate in that case.

\begin{corollary}[Selection with an economic tolerance]\label{cor:r18selection}
Let $c_i\geq0$ be observed accounting costs and $\tau\geq0$. Define
$\mathcal C_\tau=\{i:\max_jD_{ij}\leq\tau\}$.
When this set is nonempty, choose a deterministic tie-broken minimizer
$\widehat i$ of $c_i$ over $\mathcal C_\tau$.
Then on $\mathcal E$,
$\max_jV_j-V_{\widehat i}\leq\tau$ and
$c_{\widehat i}\leq c_k$ for every $k\in\mathcal C_\tau$.
The accuracy assertion remains valid for a tolerance chosen after inspecting
the intervals, interpreted as a simultaneous assertion for all tolerances.
A common nonnegative assessment charge preserves the cost ordering.
\end{corollary}
\begin{proof}
On $\mathcal E$ all candidate bounds hold at once. Membership of the chosen
candidate in $\mathcal C_\tau$ gives its regret bound irrespective of how
the index or tolerance was selected. The cost inequality is the definition
of the minimizer; addition of the same charge to both sides preserves it.
If $\mathcal C_\tau$ is empty, the rule returns no certified candidate.
\end{proof}

The corollary gives no assertion about $c_k$ for an uncertified candidate
with unknown true regret below $\tau$. Observed cost is not an expectation
over future machines or training runs. Selecting after the full computation
does not save the work already spent computing the competing candidates.
In particular, a retrospective least certified prefix allocation is not a
measured work-to-target stopping time. These distinctions are necessary
when the same continuation is amortized across an economic decision menu.

\subsection{Exact arithmetic and the complete input family}
The computation reads all 216 original protected intervals: eight economic
cells, three stages, five gains relative to the common reference, and four
direct NBO contrasts at each stage. The stored binary64 endpoints are
interpreted as exact rational numbers. Floyd--Warshall updates use only
rational addition, ordering, and minima. Final lower endpoints are rounded
downward and upper endpoints upward. The reference plus fifteen method-stage
candidates are retained in every cell, including costly and dominated
candidates. All means use the same sixteen streams; no stream is selected
on the basis of the analysis. The calculation adds no observations and
inherits the original $.018$ payoff-family failure probability.
